% META: {"id":"1SPE-VARALEA-CO-040","chapitre":"1SPE-VARIABLES-ALEATOIRES","type_objet":"corrige","exercice_ref":"1SPE-VARALEA-EX-040","capacites_codes":[],"extension_codes":["X2"],"programme_alignment":"OPTIONAL_EXTENSION","extension_label":"Approfondissement — Vers la Terminale","sources_inspiration":[],"status":"approved","origin":"generated","status_history":["generated","audited","accepted","approved"],"release_acceptance":"RELEASE_OWNER_BATCH_ACCEPTANCE","acceptance_closure_digest":"sha256:9b3ccf9a81c5520fbb7e03b7b2d3e7bf2057a02834b6d908bf3be5f49f3b553f","acceptance_receipt_digest":"sha256:4fad86f0282e0fa4e0419cca1ad6379ae7af5a280c04a4a90880f90a1c044242"}
\begin{center}\textbf{Approfondissement — Vers la Terminale}\end{center}
% BEGIN-VERIFY
% from sympy import Rational, sqrt
% n, p = 50, Rational(1,10)
% assert n*p == 5
% assert n*p*(1-p) == Rational(9,2)
% END-VERIFY
\begin{corrige}{1SPE-VARALEA-EX-040}

\textbf{1.} $X \sim \mathcal{B}(50, 0{,}1)$.

\textbf{2.} $E(X) = 50 \times 0{,}1 = 5$. $V(X) = 50 \times 0{,}1 \times 0{,}9 = 4{,}5$. $\sigma(X) = \sqrt{4{,}5} \approx 2{,}12$.

\textbf{3.} En moyenne, on trouve $5$ pièces défectueuses.

\end{corrige}
